\documentclass[11pt]{article}
\usepackage[margin=1in]{geometry}
\usepackage{amsmath,amssymb,amsthm,hyperref}
\newtheorem{theorem}{Theorem}
\newtheorem{lemma}{Lemma}
\newcommand{\conv}{\operatorname{conv}}
\newcommand{\Pol}{\operatorname{Pol}}
\title{Root localization for a derivative convolution}
\author{Principia Math In Progress}
\date{11 September 2026}
\begin{document}
\maketitle
\begin{abstract}
For complex polynomials $P,Q$ of the same degree $p$, every zero of
$\sum_{k=0}^p P^{(k)}Q^{(p-k)}$ lies in half the Minkowski sum of their
root convex hulls. This proves the derivative-convolution conjecture in
Lachaume's paper, including its complex Gauss--Lucas formulation.
The proof identifies the sum with a translated polarization and applies
the classical Grace--Walsh--Szeg\H{o} theorem.
\end{abstract}

\section*{Task 1: Resolve MathDB problem 338079}
Let $p\geq1$ and let $P,Q\in\mathbb C[X]$ both have degree exactly $p$.
Write their leading coefficients as $a,b\ne0$, list their roots with
multiplicity as $\alpha_1,\ldots,\alpha_p$ and $\beta_1,\ldots,\beta_p$,
and set
\[
 S_{P,Q}(x)=\sum_{k=0}^pP^{(k)}(x)Q^{(p-k)}(x),\qquad
 K_P=\conv\{\alpha_i\},\quad K_Q=\conv\{\beta_i\}.
\]
Theorem 4.9 of \cite{lachaume} states the real-rooted formulation of
Conjecture 4: $S_{P,Q}$ is real-rooted whenever $P,Q$ are real-rooted.
Conjecture 5 on the same page asks whether, for complex $P,Q$, every root
belongs to $\conv(K_P\cup K_Q)$.

\begin{theorem}\label{thm:main}
The polynomial $S_{P,Q}$ has degree $p$, leading coefficient $2^pp!ab$,
and all its roots belong to
\[
 \tfrac12(K_P+K_Q)
 :=\{(u+v)/2:u\in K_P,\ v\in K_Q\}
 \ \subseteq\ \conv(K_P\cup K_Q).
\]
In particular, both conjectures hold, with repeated roots permitted.
If the roots are real, every root of $S_{P,Q}$ belongs to
\[
 \left[\frac{\min_i\alpha_i+\min_i\beta_i}{2},
       \frac{\max_i\alpha_i+\max_i\beta_i}{2}\right].
\]
\end{theorem}

\subsection*{Two polynomial identities}
For a polynomial $Q(z)=\sum_{j=0}^p q_jz^j$, its polarization is
\[
 F(z_1,\ldots,z_p)=\Pol Q(z_1,\ldots,z_p)
   :=\sum_{j=0}^p\frac{q_j}{\binom pj}e_j(z_1,\ldots,z_p),
\]
where $e_j$ is the $j$th elementary symmetric polynomial and $e_0=1$.
Thus $F$ is symmetric, is affine in each variable separately, and
$F(z,\ldots,z)=Q(z)$.

\begin{lemma}\label{lem:identity}
For all $x\in\mathbb C$,
\begin{equation}\label{eq:identity}
 S_{P,Q}(x)=p!a\,F(2x-\alpha_1,\ldots,2x-\alpha_p).
\end{equation}
\end{lemma}
\begin{proof}
Define the bivariate polynomial
\[
 T(u,v)=\sum_{k=0}^pP^{(k)}(u)Q^{(p-k)}(v).
\]
Reindexing the differentiated sums gives $\partial_uT=\partial_vT$:
the potential extra terms contain $P^{(p+1)}$ or $Q^{(p+1)}$ and vanish.
Consequently, for fixed $u,v$, the polynomial
$t\mapsto T(t,u+v-t)$ has derivative zero. In particular,
\begin{equation}\label{eq:translation}
 S_{P,Q}(x)=T(x,x)=T(0,2x)
    =\sum_{k=0}^pP^{(p-k)}(0)Q^{(k)}(2x).
\end{equation}

Let $(p)_k=p!/(p-k)!$, with $(p)_0=1$.
For each subset $A\subseteq\{1,\ldots,p\}$ of size $k$, symmetry and
the diagonal identity give
\[
 (\partial_A F)(z,\ldots,z)=\frac{Q^{(k)}(z)}{(p)_k}.
\]
Indeed, differentiating the diagonal $k$ times sums the mixed derivatives
over the $(p)_k$ ordered lists of distinct indices; all terms with a
repeated index vanish because $F$ is affine in that variable.
The exact multiaffine Taylor expansion now yields
\[
 F(z-\alpha_1,\ldots,z-\alpha_p)
   =\sum_{k=0}^p(-1)^ke_k(\alpha_1,\ldots,\alpha_p)
                      \frac{Q^{(k)}(z)}{(p)_k}.
\]
On the other hand, expanding $P(z)=a\prod_i(z-\alpha_i)$ gives
\[
 \frac{P^{(p-k)}(0)}{p!a}
   =\frac{(-1)^ke_k(\alpha_1,\ldots,\alpha_p)}{(p)_k}.
\]
Substitution into \eqref{eq:translation} proves \eqref{eq:identity}.
\end{proof}

\subsection*{Localization by half-planes}
We use the following classical form of the Grace--Walsh--Szeg\H{o}
coincidence theorem \cite[Lemma 6.13, p.~106]{pemantle}: if a symmetric
multiaffine complex polynomial $F$ is evaluated at points
$z_1,\ldots,z_p$ in an open half-plane $H$, then its value equals
$F(\zeta,\ldots,\zeta)$ for some $\zeta\in H$.
An open half-plane is a convex circular region, so it satisfies the
hypothesis of the cited theorem. No reality, sign, or simplicity
condition on the coefficients or roots is required.

\begin{proof}[Proof of Theorem \ref{thm:main}]
Suppose that $x\notin\tfrac12(K_P+K_Q)$. This set is compact and
convex in the real plane $\mathbb C$. Strict separation supplies a
nonzero complex number $\eta$ such that
\[
 2\Re(\eta x)>A+B,\qquad
 A=\max_i\Re(\eta\alpha_i),\quad
 B=\max_i\Re(\eta\beta_i).
\]
Here the support function of a Minkowski sum is the sum of the two
support functions, as follows directly by taking the two maxima
independently. For every $i$,
\[
 \Re\bigl(\eta(2x-\alpha_i)\bigr)>B.
\]
Thus all $2x-\alpha_i$ lie in the open half-plane
$H=\{z:\Re(\eta z)>B\}$, while $Q$ has no root in $H$.
The coincidence theorem gives a $\zeta\in H$ for which
\[
 F(2x-\alpha_1,\ldots,2x-\alpha_p)=F(\zeta,\ldots,\zeta)
   =Q(\zeta)\ne0.
\]
Lemma \ref{lem:identity} therefore shows $S_{P,Q}(x)\ne0$.
This proves the asserted localization.

Finally, the degree-$p$ term of $F$ is $b\prod_i z_i$. Equation
\eqref{eq:identity} gives the nonzero leading coefficient $2^pp!ab$.
For real-rooted $P,Q$, the two convex hulls are their real root
intervals, whose half-sum is exactly the interval displayed in the
theorem. This also proves real-rootedness.
\end{proof}

The enclosure cannot be uniformly reduced even in the one-point case:
if $P(x)=a(x-\alpha)^p$ and $Q(x)=b(x-\beta)^p$, then
\[
 S_{P,Q}(x)=p!ab(2x-\alpha-\beta)^p.
\]
The localization is an application of classical polarization theory;
no novelty is claimed for the coincidence theorem or for the
underlying polynomial convolution.

\subsection*{Scope and status}
This note fully answers MathDB 338079. The complex theorem also answers
the companion entry 338080, which receives no separate campaign credit.
The source and the author's thesis retain the exact two conjectures.
The author's later gravity paper \cite[Theorem 2.7(d)]{gravity} claims a
weaker concavity consequence using Walsh's theorem; it does not state
a solution of these derivative-convolution conjectures. Dated live
records and the bounded search for later claims are documented in the
accompanying literature audit. The proof has been checked locally;
external peer review is not claimed.

\section*{Task 2: Determine boundary-root multiplicities}
Fix a supporting line of $\tfrac12(K_P+K_Q)$. Determine exactly when
$S_{P,Q}$ has a root on that line and express its multiplicity in terms
of the roots of $P,Q$ on the corresponding supporting lines. Include
coincident roots and line segments, and compare the result with the
multiplicity theorems for classical polynomial convolution.

\section*{Task 3: Classify admissible derivative weights}
For fixed $p$, characterize the real vectors $(c_0,\ldots,c_p)$ for which
every nonzero polynomial
\[
 \sum_{k=0}^p c_kP^{(k)}Q^{(p-k)}
\]
has its roots in $\conv(K_P\cup K_Q)$ for every pair of complex
degree-$p$ polynomials. Obtain an explicit criterion, accounting for
zero outputs and arbitrary nonzero input leading coefficients.

\section*{Task 4: Extend the identity to several inputs}
For $r\geq3$ degree-$p$ polynomials, analyze the sum of
$\prod_{j=1}^rP_j^{(k_j)}$ over $0\leq k_j\leq p$ with
$\sum_j k_j=(r-1)p$. Derive its exact normalization and determine its
root enclosure in terms of the average of the $r$ root convex hulls.
Give the relation to iterated binary convolution explicitly.

\section*{Task 5: Quantify movement of the output roots}
For real-rooted inputs whose normalized convolution has simple roots,
derive bounds for the movement of each output root when the input roots
are perturbed. State the dependence on input diameter, degree, and
the smallest output-root separation, and identify the deterioration
as an output root becomes multiple.

\begin{thebibliography}{9}
\bibitem{lachaume}
X. Lachaume, \emph{On the concavity of a sum of elementary symmetric
polynomials}, arXiv:1712.10327v1 (2017), Theorem 4.9 and Conjecture 5,
p.~17. \url{https://arxiv.org/abs/1712.10327}.
\bibitem{pemantle}
R. Pemantle, \emph{Hyperbolicity and stable polynomials in combinatorics
and probability}, Current Developments in Mathematics 2011,
57--123, Lemma 6.13, p.~106.
\url{https://intlpress.com/site/pub/files/_fulltext/journals/cdm/2011/2011/0001/CDM-2011-2011-0001-a002.pdf}.
\bibitem{gravity}
X. Lachaume, \emph{The constraint equations of Lovelock gravity theories:
a new $\sigma_k$-Yamabe problem}, arXiv:1712.04528v2 (2018),
Theorem 2.7(d), p.~8. \url{https://arxiv.org/abs/1712.04528}.
\end{thebibliography}
\end{document}
